Understanding how deep neural networks learn useful internal representations from data remains a central open problem in the theory of deep learning. We introduce \emph{Neural Low-Degree Filtering} (Neural LoFi), a stylized limit of gradient-based training in which hierarchical feature learning becomes an explicit iterative spectral procedure. In this limit, the dynamics at each layer decouple: given the current representation, the next layer selects directions with maximal accessible low-degree correlation to the label. This yields a tractable surrogate mechanism for deep learning, together with a natural kernel-space interpretation. Neural LoFi provides a mathematically explicit framework for studying multi-layer feature learning beyond the lazy regime. It predicts how representations are selected layer by layer, explains how emergence of concepts arises with given sample complexity,
and gives a concrete mechanism by which depth progressively constructs new features from old ones through low-degree compositionality. We complement the theory with mechanistic experiments on fully connected and convolutional architectures, showing that Neural LoFi improves over lazy random-feature baselines, recovers meaningful structured filters, and predicts representations aligned with early gradient-descent feature discovery with real datasets.
